\begin{ouverture}[courbe]{12}{Statistiques à deux variables}{Statistique et esprit critique}{Nuage · Régression · Ajustement · Causalité}
\objectif{Représenter un nuage et calculer son point moyen.}
\objectif{Déterminer et utiliser un ajustement affine.}
\objectif{Linéariser un modèle exponentiel.}
\objectif{Distinguer corrélation, prévision et causalité.}
\prerequis{Moyenne, variance et fonction affine.}
\end{ouverture}
\essentiel

\essSection{1}{Deux mesures sur les mêmes individus}
Une série à deux variables est une liste de couples $(x_i,y_i)$ observés sur les mêmes individus ou aux mêmes dates. Le nuage de points conserve l'association entre les deux mesures.
Le point moyen est $G(\bar x,\bar y)$, avec $\bar x=\sum x_i/n$ et $\bar y=\sum y_i/n$.
Un nuage peut suggérer une tendance affine, exponentielle ou autre. Le choix d'un modèle doit être discuté à partir du contexte et des résidus, pas seulement d'un indicateur.
\figurelivre{c12-regression.pdf}{Données fictives $(1,2),(2,3),(3,5),(4,6)$. La droite des moindres carrés est $y=1{,}4x+0{,}5$ et passe par $G(2{,}5,4)$. Les segments verticaux sont les résidus.}
\essSection{2}{La droite des moindres carrés}
\begin{definition}
L'ajustement de $y$ en $x$ cherche la droite $y=ax+b$ minimisant
$S(a,b)=\sum_{i=1}^n(y_i-ax_i-b)^2$.
Les distances mesurées sont verticales. L'ajustement de $x$ en $y$ est un autre problème.
\end{definition}
Avec $V_x=\sum(x_i-\bar x)^2/n>0$ et
$\operatorname{Cov}(x,y)=\sum(x_i-\bar x)(y_i-\bar y)/n$,
\[a=\frac{\operatorname{Cov}(x,y)}{V_x},\qquad b=\bar y-a\bar x.\]
La calculatrice ou un tableur fournit ces coefficients. On doit annoncer la variable expliquée et conserver assez de chiffres dans les calculs.
\begin{demonstration}[Pourquoi ces coefficients ?]
Poser $X_i=x_i-\bar x$ et $Y_i=y_i-\bar y$. Comme $\sum X_i=\sum Y_i=0$,
$S(a,b)=\sum(Y_i-aX_i)^2+n(\bar y-a\bar x-b)^2$.
Pour $a$ fixé, le deuxième terme est minimal si $b=\bar y-a\bar x$.
Le premier vaut $n(V_y-2a\operatorname{Cov}+a^2V_x)$, trinôme minimal en $a=\operatorname{Cov}/V_x$.
\end{demonstration}
\essSection{3}{Le coefficient de corrélation}
Si les deux variances sont non nulles,
$r=\operatorname{Cov}(x,y)/\sqrt{V_xV_y}$, avec $-1\leq r\leq1$.
Une valeur proche de $1$ ou de $-1$ indique une forte liaison affine dans les données. Un $r$ proche de zéro ne signifie pas absence de toute relation : un nuage peut suivre une parabole symétrique.
Une valeur aberrante peut fortement déplacer la droite et $r$.
\essSection{4}{Interpoler, extrapoler et transformer}
Interpoler consiste à prédire entre les valeurs observées de $x$ ; extrapoler, au-delà. Plus l'on s'éloigne du domaine observé, plus la pertinence du modèle doit être examinée.
Si $y>0$ et si un modèle $y=A\mathrm e^{kx}$ est plausible, poser $z=\ln y$ donne $z=\ln A+kx$.
Après une régression de $z$ en $x$, on retrouve $A=\mathrm e^b$, $k=a$.
Cet ajustement minimise les résidus au carré sur l'échelle logarithmique ; ce n'est pas nécessairement le même critère que les résidus sur l'échelle initiale.
\begin{exemple}
Avec $y=2\times3^x$, $\ln y=\ln2+x\ln3$. Des rapports constants entre valeurs successives suggèrent un modèle exponentiel, tandis que des différences constantes suggèrent un modèle affine.
\end{exemple}
\essSection{5}{Corrélation et causalité}
Une corrélation peut venir d'un effet causal, d'un facteur commun, d'une causalité inverse, d'une sélection des observations ou d'une coïncidence.
\figurepetite{c12-causalite.pdf}{Un facteur commun peut influencer deux variables corrélées. La corrélation entre elles ne suffit pas à établir une flèche causale directe. Schéma conceptuel fictif.}
Un protocole, la connaissance du phénomène et le contrôle des facteurs confondants sont nécessaires pour argumenter une causalité.
À l'oral, présenter les axes, le domaine des observations, le modèle, ses résidus et ses limites. Éviter de transformer une association statistique en explication certaine.

\exercices
\markright{Prendre confiance}
\begingroup\renewcommand{\otBesoinPlace}[1]{}\begin{minipage}{\linewidth}
\rubrique{Prendre confiance}
\begin{exercice}[1]{Point moyen}
Les couples sont $(1,3),(2,4),(3,8)$.
\begin{questions}
\item Calculer $\bar x$.
\item Calculer $\bar y$.
\item Donner le point moyen.
\end{questions}
\end{exercice}
\begin{exercice}[1]{Utiliser une droite}
Un ajustement donne $y=2{,}5x+7$ pour des valeurs observées de $x$ entre $2$ et $8$.
\begin{questions}
\item Prédire $y$ pour $x=6$.
\item Prédire pour $x=20$.
\item Qualifier les deux prédictions.
\end{questions}
\end{exercice}
\begin{exercice}[1]{Résidu}
La droite ajustée est $y=3x+1$. Un point observé est $(2,8)$.
\begin{questions}
\item Calculer la valeur prédite.
\item Calculer le résidu observé moins prédit.
\item Donner sa contribution à la somme des carrés.
\end{questions}
\end{exercice}
\begin{exercice}[1]{Lire une corrélation}

\begin{questions}
\item Que suggère $r=-0{,}98$ ?
\item Que ne permet-il pas d’affirmer ?
\item Pourquoi $r=0$ n’exclut-il pas une relation non affine ?
\end{questions}
\end{exercice}
\end{minipage}\endgroup
\clearpage
\markright{Appliquer et relier}
\begingroup\renewcommand{\otBesoinPlace}[1]{}\begin{minipage}{\linewidth}
\rubrique{Appliquer et relier}
\begin{exercice}[2]{Calculer une régression}
On observe $(1,2),(2,3),(3,5),(4,6)$.
\begin{questions}
\item Calculer les moyennes, $V_x$ et la covariance.
\item Déterminer $a,b$.
\item Calculer $r$, sachant $V_y=2{,}5$.
\end{questions}
\end{exercice}
\begin{exercice}[2]{Exponentielle par transformation}
Des données suivent exactement $y=5\times2^x$.
\begin{questions}
\item Exprimer $z=\ln y$.
\item Donner la pente et l’ordonnée à l’origine dans le repère $(x,z)$.
\item Retrouver $y$ à partir de ces coefficients.
\end{questions}
\end{exercice}
\begin{exercice}[2]{Interpolation et seuil}
Un modèle fictif est $y=4\mathrm e^{0{,}3t}$, ajusté pour $0\leq t\leq5$.
\begin{questions}
\item Prédire la valeur en $t=3$.
\item Résoudre $y=20$.
\item Discuter le résultat par rapport aux données.
\end{questions}
\end{exercice}
\begin{exercice}[2]{Moindres carrés}
Pour trois points, une droite $D_1$ donne les résidus $1,-1,0$ ; $D_2$ donne $0,0,2$.
\begin{questions}
\item Calculer les deux sommes de carrés.
\item Quelle droite est préférable selon ce critère parmi les deux ?
\item Peut-on affirmer qu’il s’agit de la meilleure de toutes les droites ?
\end{questions}
\end{exercice}
\end{minipage}\endgroup
\clearpage
\markright{Chercher et argumenter}
\begingroup\renewcommand{\otBesoinPlace}[1]{}\begin{minipage}{\linewidth}
\rubrique{Chercher et argumenter}
\begin{exercice}[3]{Corrélation nulle et parabole}
On considère $(-1,1),(0,0),(1,1)$.
\begin{questions}
\item Calculer le point moyen.
\item Montrer que la covariance vaut zéro.
\item Quelle relation exacte existe pourtant ?
\end{questions}
\end{exercice}
\begin{exercice}[3]{Une variable confondante}
Dans des données fictives, les ventes de boissons fraîches et les baignades augmentent ensemble.
\begin{questions}
\item Proposer un facteur commun plausible.
\item Peut-on conclure que les achats provoquent les baignades ?
\item Que faudrait-il examiner pour discuter un lien causal ?
\end{questions}
\end{exercice}
\begin{exercice}[3]{Changer les unités}
Une régression relie $x$ en heures à $y$ en kilomètres : $y=60x+5$.
\begin{questions}
\item Écrire le modèle avec le temps $u$ en minutes.
\item La valeur prédite pour le même instant change-t-elle ?
\item Le coefficient de corrélation change-t-il sous ce changement d’unité positif ?
\end{questions}
\end{exercice}
\begin{exercice}[3]{Pourquoi la droite passe par $G$}

\begin{questions}
\item Pour $a$ fixé, écrire $S(a,b)$ en séparant le terme contenant $\bar y-a\bar x-b$.
\item En déduire le meilleur $b$.
\item Que se passe-t-il si tous les $x_i$ sont égaux ?
\end{questions}
\end{exercice}
\end{minipage}\endgroup
